\documentclass[10pt,a4paper]{amsart}
\usepackage[margin=19mm]{geometry}
\usepackage{amsmath,amssymb,mathtools}
\usepackage[scr=rsfs,cal=euler]{mathalfa}
\usepackage{xfrac}
\usepackage[unicode,hidelinks,bookmarks=false]{hyperref}
\newcommand{\ie}{i.e.\ }
\newcommand{\cf}{cf.\ }
\newcommand{\resp}{respectively\ }
\newcommand{\Subsection}{\subsection}
\providecommand{\nobreakdash}{\nobreak\hbox{-}}
\setlength{\emergencystretch}{2em}
\multlinegap0pt
\usepackage{amssymb}
\usepackage{mathtools}
\newtheorem{lemma}{Lemma}
\newcommand{\e}{\mathrm e}
\newcommand{\ind}{\mathbf 1}
\title{Unbounded logarithmic limsup in Erd\H{o}s Problem 684\\ via shifted carry scheduling}
\author{Ji Ho Bae}
\date{Source: arXiv:2604.23784v3 (CC BY 4.0)}
\begin{document}
\maketitle

\noindent\emph{Source and licence.} Unbounded logarithmic limsup in Erd\H{o}s Problem 684\\ via shifted carry scheduling. Version arXiv:2604.23784v3 (CC BY 4.0), distributed by arXiv under the Creative Commons Attribution 4.0 International licence, http://creativecommons.org/licenses/by/4.0/, as recorded in the arXiv licence metadata for this exact version and shown on the abstract page https://arxiv.org/abs/2604.23784v3. Full licence notice: \texttt{licenses/arxiv-2604.23784-CC-BY-4.0.txt}.\par\smallskip

\noindent\textit{Editorial scope.} Counted unit: the article's lemma lem:tail (source0.tex lines 388--398). The article's complete original proof of it is delivered with it (source0.tex lines 400--579, one \texttt{proof} environment of 180 lines). No mathematics is written, repaired or re-derived: the statement and the proof are the article's own lines, in the article's own notation, and no retained line is substituted or re-flowed.  Retained uncounted as the source-local context the proof needs: lines 236--240; lines 261--264.  Nothing was omitted from any retained span, so the package carries no omission record.  No retained line contains a citation, so the wrapper carries no bibliography and no bibliography entry.  No retained line contains a figure environment, an image-inclusion command or drawing code, so no image is delivered and the package declares no asset dependency.  Editorial formatting only, in addition: an AMS \texttt{amsart} preamble that restates the article's own declarations and macros used by the retained lines, namely the environments \texttt{lemma} on separate counters, together with the standard packages those lines need.  The retained environments print in this document as Lemma 1 in order of appearance, whereas the article prints the same items with the numbering of its own full document.  The complete native source of the article as distributed in the arXiv e-print for this exact version is bundled unchanged as \texttt{sources/199190/source0.tex}.
\begin{equation}\label{eq:psi-theta}
 \psi(x)-\vartheta(x)=O(\sqrt{x}),
 \qquad
 \psi(x)=x+o(x).
\end{equation}

\begin{equation}\label{eq:param-basic}
 0<\lambda<1,\qquad \beta>0,\qquad
 0<c<\lambda\beta,\qquad \sigma>2\beta,
\end{equation}

\begin{lemma}[High-power tail]\label{lem:tail}
For every fixed parameter tuple satisfying \eqref{eq:param-basic},
\begin{equation}\label{eq:tail-bound}
 \#\{1\leq t\leq J:Z(t)\geq\beta M\}
 \leq
 2J\exp\!\left(-\lambda\beta M
       +O_\lambda\bigl(K^{(1+\lambda)/2}\bigr)\right).
\end{equation}
In particular, because \(\lambda<1\) is fixed and
\(K=M\log^{1+o(1)}M\), the error in the exponential is \(o(M)\).
\end{lemma}

\begin{proof}
We divide the proof into four steps.

\medskip
\noindent\emph{Step 1: nested towers.}
For \(p\leq K\), define
\begin{equation}\label{eq:b-a}
 b_p=\nu_p(L),
 \qquad
 a_p=\min\{a\geq2:p^a>K\}.
\end{equation}
Since \(p^{b_p}\leq M<K<p^{a_p}\), one has \(a_p>b_p\).
For \(a\geq a_p\), the events
\[
 E_{p,a}(t):=\{[n_t]_{p^a}<K\}
\]
are nested downwards in \(a\).  Indeed, if \(E_{p,a+1}(t)\) holds, its
residue is less than \(K<p^a\), so reduction modulo \(p^a\) leaves the
same residue and \(E_{p,a}(t)\) holds.  Let \(\ell_p(t)\) be the length of
the initial true segment beginning at \(a_p\), with length zero if the first
event fails.  The preceding finiteness observation gives
\begin{equation}\label{eq:Z-depth}
 Z(t)=\sum_{p\leq K}\ell_p(t)\log p.
\end{equation}

\medskip
\noindent\emph{Step 2: congruence classes for a depth profile.}
Put \(g_p=p^{b_p}\).  If \(\ell_p(t)\geq r\geq1\), then for some integer
\(0\leq s<K\),
\begin{equation}\label{eq:depth-congruence}
 tL\equiv h+1+s\pmod {p^{a_p+r-1}}.
\end{equation}
Solvability forces \(g_p\mid h+1+s\).  There are at most
\(K/g_p+1\leq2K/g_p\) possible values of \(s\), since
\(g_p\leq M<K\).  After division by \(g_p\), the integer \(L/g_p\) is a
unit modulo
\begin{equation}\label{eq:Qpr}
 Q_p(r)=p^{a_p+r-1-b_p}.
\end{equation}
Consequently, the condition \(\ell_p(t)\geq r\) restricts \(t\) to at
most \(2K/g_p\) residue classes modulo \(Q_p(r)\).

For a finitely supported profile
\(\mathbf r=(r_p)_{p\leq K}\) of nonnegative integers, define
\begin{equation}\label{eq:RQ}
 R(\mathbf r)=\prod_{r_p>0}\frac{2K}{g_p},
 \qquad
 Q(\mathbf r)=\prod_{r_p>0}Q_p(r_p),
 \qquad
 W(\mathbf r)=\sum_{p\leq K}r_p\log p.
\end{equation}
The moduli \(Q_p(r_p)\) have distinct prime bases.  The Chinese remainder
theorem and elementary interval counting therefore give
\begin{equation}\label{eq:crt-plus-one}
 \#\{1\leq t\leq J:\ell_p(t)\geq r_p\text{ for every }p\}
 \leq R(\mathbf r)\left(\frac{J}{Q(\mathbf r)}+1\right).
\end{equation}

\medskip
\noindent\emph{Step 3: a short-modulus witness for every bad multiplier.}
Suppose that \(Z(t)\geq\beta M\).  Starting from the zero profile, add
one unit at a time to coordinates without exceeding
\(\ell_p(t)\), and stop when the weight first reaches \(\beta M\).
Since an increment has weight at most \(\log K\), the resulting prefix
profile satisfies
\begin{equation}\label{eq:minimal-prefix}
 \beta M\leq W(\mathbf r)<\beta M+\log K.
\end{equation}
Let \(\mathcal W_M\) be the finite set of all nonnegative profiles satisfying
\eqref{eq:minimal-prefix}.  Every multiplier with \(Z(t)\geq\beta M\)
belongs to the depth event associated with at least one profile in
\(\mathcal W_M\).  We now prove that every profile in \(\mathcal W_M\) has
reduced CRT modulus below the search length \(J\).

From \eqref{eq:Qpr} and \eqref{eq:RQ},
\begin{equation}\label{eq:logQ}
 \log Q(\mathbf r)
 =W(\mathbf r)+
  \sum_{r_p>0}(a_p-1-b_p)\log p.
\end{equation}
For \(p>M\), the relation \(K<M^2<p^2\) gives
\(a_p=2\) and \(b_p=0\); hence that prime's summand in the second term
is \(\log p\leq r_p\log p\).  For \(p\leq M\), one has
\[
 a_p-1=\max\{a:p^a\leq K\},
 \qquad
 b_p=\max\{b:p^b\leq M\}.
\]
Thus, on summing even over all primes rather than merely the support of
\(\mathbf r\), cancellation of the first-power terms gives the exact identity
\begin{equation}\label{eq:overhead-identity}
 \sum_{p\leq M}(a_p-1-b_p)\log p
 =\{\psi(K)-\vartheta(K)\}
  -\{\psi(M)-\vartheta(M)\}.
\end{equation}
Here \(K<M^2\) ensures that every prime having a square at most \(K\) is
at most \(M\).  The expression in \eqref{eq:overhead-identity} is
nonnegative and is \(O(\sqrt K)=o(M)\) by
\eqref{eq:psi-theta}.  It follows from
\eqref{eq:minimal-prefix}--\eqref{eq:overhead-identity} that
\begin{equation}\label{eq:Q-before-plus-one}
 \log Q(\mathbf r)
 \leq2W(\mathbf r)+O(\sqrt K)
 \leq2\beta M+o(M).
\end{equation}
Put \(\varepsilon_0=(\sigma-2\beta)/3>0\).  Since
\(\log J=\sigma M+o(1)\), for all sufficiently large \(M\) we have
\[
 \log Q(\mathbf r)
 \leq(2\beta+\varepsilon_0)M
 <(\sigma-\varepsilon_0)M
 <\log J.
\]
Thus
\begin{equation}\label{eq:Q-less-J}
 Q(\mathbf r)<J
\end{equation}
for every \(\mathbf r\in\mathcal W_M\).

For every \(\mathbf r\in\mathcal W_M\), \eqref{eq:Q-less-J} permits
\(J/Q(\mathbf r)+1\leq2J/Q(\mathbf r)\) in
\eqref{eq:crt-plus-one}.  The factors
\(g_p\) cancel exactly, and \eqref{eq:crt-plus-one} becomes
\begin{equation}\label{eq:profile-count}
 \#\{1\leq t\leq J:\ell_p(t)\geq r_p\text{ for every }p\}
 \leq
 2J\prod_{r_p>0}\frac{2K}{p^{a_p+r_p-1}}.
\end{equation}
This also covers the original tail containing individual prime powers larger
than \(J\): only the reduced joint modulus of the extracted prefix is used,
and \eqref{eq:Q-less-J} places that modulus in the legitimate range.

\medskip
\noindent\emph{Step 4: exponential weighting.}
Take the union in \eqref{eq:profile-count} over
\(\mathbf r\in\mathcal W_M\).  Apply exponential weighting in the form
\(\ind\{W\geq\beta M\}\leq
\e^{-\lambda\beta M}\e^{\lambda W}\).  Enlarging the resulting
\emph{nonnegative numerical sum} to all profiles gives
\begin{align}
 \#\{1\leq t\leq J:Z(t)\geq\beta M\}
 &\leq2J\e^{-\lambda\beta M}
 \prod_{p\leq K}
 \left(1+\sum_{r\geq1}
   \frac{2Kp^{\lambda r}}{p^{a_p+r-1}}\right)\notag\\
 &=2J\e^{-\lambda\beta M}
 \prod_{p\leq K}
 \left(1+
   \frac{2K}{p^{a_p-1}(p^{1-\lambda}-1)}\right).
 \label{eq:rankin-product}
\end{align}
The last enlargement is purely numerical: the CRT estimate has already been
applied to every profile in \(\mathcal W_M\).

For fixed \(\lambda<1\), one has
\(p^{1-\lambda}-1\gg_\lambda p^{1-\lambda}\) uniformly for \(p\geq2\).
Consequently, the logarithm of the product in
\eqref{eq:rankin-product} is at most a constant depending on \(\lambda\)
times
\[
 \sum_{p\leq\sqrt K}p^\lambda
 +K\sum_{p>\sqrt K}p^{-(2-\lambda)}.
\]
Indeed, for \(p\leq\sqrt K\), the minimality of \(a_p\) gives
\(p^{a_p-1}\leq K<p^{a_p}\), whereas for \(p>\sqrt K\) one has
\(a_p=2\).  Enlarging the prime sums to integer sums shows that the last
display is
\begin{equation}\label{eq:euler-error}
 O_\lambda\bigl(K^{(1+\lambda)/2}\bigr).
\end{equation}
Using \(\log(1+x)\leq x\) proves \eqref{eq:tail-bound}.
Finally,
\[
 \frac{K^{(1+\lambda)/2}}M
 =(1+o(1))M^{-(1-\lambda)/2}
  \left(c\frac{\log M}{\log\log M}\right)^{(1+\lambda)/2}
 =o(1)
\]
for every fixed \(\lambda<1\).
\end{proof}

\end{document}
